\section{What this extension adds to the existing programme}
\label{sec:source-audit}

\subsection{Audit boundary and proof interface}

The source comparison is pinned to ProveIt commit
\nolinkurl{3ee850b4c2d650cf01263823fd5415963747c5d0}.
It uses the current unification plan and derivation inventory, the proof-packet
index, and the particular analytic and foundational chapters cited below.
It is a focused audit of the relevant interfaces, not a verification of every
research packet. The repository already contains many frameworks for bounded
certificates, finite observations, scoped witnesses, reverse precision
demands, and residual hypotheses. Merely giving those constructions another
name would not resolve a new mathematical difficulty.

The maintained interface is ordinary intuitionistic primitive recursive
arithmetic: intuitionistic first-order logic with equality, defining equations
for primitive recursive functions, successor axioms, and quantifier-free
induction. In the Lean development the public object judgment is
\(\Gamma\vdash[\texttt{PRAAx}]\varphi\).
The current unification plan expressly distinguishes an object derivation
from an ambient Lean theorem about standard natural numbers, and the
derivation inventory records the precise assumptions of each checked
endpoint \cite{audit-unification,audit-progress}.
Consequently a new rational estimate, a proof-producing arithmetic lemma,
and a completed analytic application have separate statuses throughout this
report.

\subsection{Scalar contractions and polynomial ODEs are established predecessors}

Two earlier analytic constructions should be credited explicitly.
The first embedding packet proves a finite scalar contraction estimate for
a normalized polynomial \(Q(X,Z)\), with contraction factor \(1/4\);
coefficient stabilization and a geometric tail then give a simple analytic
branch. Its perturbation estimate is independent of the number of iterations.
It also constructs a complex reciprocal from a finite derivative-separation
certificate, rather than assuming a reciprocal operation on arbitrary
Cauchy names \cite{audit-scalar}.

A subsequent chapter goes further. Given a root name \(b\), an inverse
derivative name \(w\), the identities
\[
 P(0,b)=0,\qquad wP_Y(0,b)=1,
\]
and numerical magnitude bounds, it chooses rational samples and small
radii by explicit formulas. It thereby \emph{produces} a passing
simple-root contraction certificate \cite{audit-scalar-producer}.
Thus the present work does not claim to introduce scalar contraction,
quantitative inversion, or the idea of producing an analytic chart.

Similarly, the local-to-global packet already gives a detailed human
construction of solutions to rational polynomial differential systems.
Dimension and degree are numerical parameters. Finite coefficient sums
produce bounds \(M,L\); a local horizon \(T=1/(2M)\) is constructed;
Euler polygons are compared through discrete Gronwall and adjacent-mesh
estimates; an explicit reindexing gives a regular solution name.
The same chapter proves uniqueness, perturbation estimates, and polynomial
first integrals \cite{audit-ode}.
Its perturbation discussion already accommodates rounded steps or
Cauchy-named coefficients when common domain and invariance contracts are
available.

The remaining formal boundary is narrower than the absence of an ODE
argument. The maintained progress inventory describes
\texttt{Foundations/DiscreteGronwall} as checked cleared arithmetic, and
\texttt{Foundations/OdeRefinement} as checked refinement arithmetic with the
error recurrence supplied. It keeps the complete rational polygon and
regular-name frontends separate \cite{audit-progress}.
A Picard or quadrature application of the present theorem is therefore an
alternative, more modular presentation of part of this analysis and an
extension to appropriately presented integrands. It is not the first
constructive polynomial ODE theorem in the repository.

\subsection{The specific gap: numerical feedback and a produced modulus}

The distinction between a cyclic precision calculation and execution of a
cyclic mathematical system is explicit in the current sources.
The checked two-coordinate demand example iterates
\[
 (x,y)\longmapsto
   \bigl(\min(B,y+1),\,\min(C,x+1)\bigr)
\]
and derives stabilization by the sum \(B+C\).
Its module documentation identifies it as a precision-demand instance, not
certificate execution \cite{audit-demand}.
The state-adaptive realizer implementation likewise handles bounded folds
using finite candidate tables and explicit paired selector alignment;
its inventory leaves broader cyclic execution and producer scopes
separate \cite{audit-adaptive,audit-progress}.

The new sufficient condition is quantitative and numerical.
For a finite family of state components, let \(A\) be a nonnegative rational
gain matrix. Supply positive rational weights \(w_i\) and a rational
\(q<1\) such that
\[
                     Aw\le q w
\]
coordinatewise. Local implementations have prescribed error budgets and
preserve a certified domain. The weighted maximum norm then turns the
coupled system into a contraction. A finite recurrence controls
iteration error and produces a Cauchy schedule, including the accumulation
of local approximation errors.

This addresses the feedback subclass for which the gain and domain
certificates hold. It does not justify circular proofs. Every local
inequality must already have an ordinary derivation; those derivations
are composed and used in a finite iteration argument. Nor does it solve
arbitrary recursive certificate systems. The numerical contraction is
additional mathematical information, and the theorem makes its use
visible.

This also strengthens the existing summable-motion interface in a useful
direction. The FIBER analysis chapter correctly distinguishes individually
valid approximations, finite stable windows, and a global Cauchy modulus.
Its telescoping theorem consumes a primitive recursive bound on every
finite tail \cite{audit-fiber}. Here the geometric tail is derived from
the gain certificate and the local error schedule. The contribution is the
\emph{producer of the tail and of a compatible sampler}; finite
telescoping and Booleanized induction are reused foundations.

\subsection{Approximate commutation extends the earlier tower methods}

There are already two distinct tower constructions in the repository.
The coherent-methods packet presents exact restrictions
\(\pi_{m,n}\) and proves coherence when a fixed producer satisfies a
specification preserved by restriction and unique at each finite level
\cite{audit-coherent}.
The topology framework uses finite simplicial complexes and refinement
carriers. Its map-coherence datum is naturality up to a common carrier,
hence contiguity, rather than equality of selected vertices
\cite{audit-topology}.
FINER's observational descent transports supplied agreement estimates
along a finite spanning tree \cite{audit-finer}.

None of these inspected chapters supplies the following analytic comparison
from a commutation defect. Let \(T_n,T_m\) be two finite-level updates,
\(R\) a restriction of Lipschitz constant \(L_R\), and suppose
\[
 \|RT_m z-T_n Rz\|\le\delta .
\]
Assume \(T_n\) contracts by \(q<1\), and the actual returned observations
satisfy
\[
 \|x_n-T_nx_n\|\le\varepsilon_n,\qquad
 \|x_m-T_mx_m\|\le\varepsilon_m .
\]
The triangle inequality and the contraction give the finite estimate
\[
 (1-q)\,\|x_n-Rx_m\|
       \le
       \varepsilon_n+L_R\varepsilon_m+\delta .
\]
Indeed insert \(T_nx_n\), \(T_nRx_m\), and \(RT_mx_m\)
between \(x_n\) and \(Rx_m\), bound the four differences, and move the
term \(q\|x_n-Rx_m\|\) to the left. This derivation uses only
finite rational inequalities; it does not assume exact fixed points.

An explicit bound on \(\delta\), together with selected numerical
precisions and a lower bound on \(1-q\), therefore converts finite-level
solvers into coherent finite observations. This is the precise new
connection with projective and spatial discretization.
The arithmetic target quantifies over finite level, observed coordinate,
iteration count, and demanded precision. An external interpretation as an
infinite sequence or inverse-limit object is justified by these contracts;
no function quantifier or compactness axiom is inserted into iPRA.

The denominator \(1-q\) is substantive. If the contraction factors
approach one, refinement defects must decrease sufficiently rapidly to
compensate, and the required schedule must be given explicitly.
Neither exact restriction coherence nor a useful limiting modulus follows
from small defects alone.

\subsection{Which further difficulties the theorem addresses}

The same finite comparison supports three reusable constructions.
A Neumann identity for finite matrix sums produces an inverse name and
a residual-to-solution error estimate whenever a rational
preconditioner has a certified contraction defect.
A multidimensional implicit system can use
\(T(x)=x-BF(x)\), with a finite box-invariance test and a bound on
\(I-BJ(x)\), to produce simultaneous root coordinates.
A finite quadrature implementation of a Picard map can use the same
error recurrence to control its discretization and evaluation errors.
Each application must derive its local bounds from its actual numerical
presentation. A claim that an inverse, derivative gap, or quadrature
modulus merely exists does not discharge that work.

These are extensions of the existing local scalar and ODE constructions
to coupled and approximately commuting computations.
They do not remove the need for separation where a problem asks for
a discontinuous label, or the need for a domain invariant on a longer
time interval. In particular, singular branch selection and global
existence beyond the certified region are not consequences of a local
contraction theorem.

\subsection{Diophantine and ineffective-constant boundaries remain distinct}

The Faltings packet distinguishes a positive height/list theorem from
a cardinality statement: a complete list and a height bound are equivalent
for its concrete point coding, but a bound on the number of points does not
tell an enumeration when to stop \cite{audit-faltings}.
This is not an assertion that every possible arithmetization of finiteness
must first produce a height bound.

The general-route chapter of the Siegel packet makes that distinction
especially clear. It proposes a finite-tuple covering theorem followed
by a zero-count theorem to rule out lists of too many distinct points.
That route could establish a cardinality conclusion without locating
all points. However, the finite approximation-cover theorem, the required
geometric presentations, and the cover-and-lift estimates are not supplied
by the finite counting combinator \cite{audit-siegel}.
The new contraction machinery can help evaluate local analytic data once
the needed quantitative contracts have been proved. It does not produce
these global arithmetic contracts.

The prime-distribution packet similarly isolates an effective
Siegel-type input needed for its requested uniform bounds, while proving
finite ingredients such as a weighted Vaughan identity and a rational Gram
certificate \cite{audit-primes}. Explicit local error budgets cannot turn
an ineffective existence proof into numerical constants.

Accordingly the resolution is partial but substantive: a broad class of
coupled and discretized constructions gains an actual primitive recursive
solver, a produced modulus, and a finite arithmetic proof route.
General Faltings, general Siegel, unrestricted branch choice, and the
missing effective prime-distribution inputs remain separate research
problems. Their unresolved status is not a proof of independence or
unprovability in iPRA.

